\[ %%%%%%%%%%%%%%%%%%%%%%%%%%%%%%%%%%%%%%%%%%%%%%%%%%%%%%%%%%%%%%%%%%%%%%%%%%%%%%%% \subsection{Scientific Challenges} %%%%%%%%%%%%%%%%%%%%%%%%%%%%%%%%%%%%%%%%%%%%%%%%%%%%%%%%%%%%%%%%%%%%%%%%%%%%%%%% Despite remarkable progress achieved by Scientific Machine Learning during the past decade, several fundamental challenges continue to limit the practical applicability of learning-based surrogate models for scientific computing. These challenges are not merely implementation issues but arise from intrinsic differences between statistical learning and physical system evolution. The first challenge concerns long-horizon stability. Most existing neural architectures are optimized using one-step prediction objectives. During inference, however, predicted states become recursive inputs for subsequent predictions. Consequently, approximation errors introduced at individual time steps accumulate over time, often resulting in exponentially increasing trajectory divergence. This phenomenon is particularly pronounced in chaotic dynamical systems where neighboring trajectories separate rapidly due to positive Lyapunov exponents. The second challenge involves preservation of physical consistency. Numerical simulations of physical systems satisfy numerous invariants, including conservation of mass, momentum, energy, charge, entropy inequalities, geometric constraints, and problem-specific constitutive relations. Conventional deep learning architectures possess no intrinsic mechanism guaranteeing that these properties remain satisfied throughout the prediction process. Even when global prediction errors remain relatively small, local violations of conservation principles may accumulate and ultimately destabilize the entire simulation. A third challenge concerns heterogeneous spatial representations. Many real-world engineering systems are naturally represented using unstructured meshes, particle discretizations, adaptive finite-volume methods, or hybrid computational domains whose topology evolves dynamically during simulation. Such representations significantly complicate the application of convolution-based neural architectures while simultaneously motivating graph-based computational models. The fourth challenge is computational scalability. High-fidelity scientific simulations frequently involve millions of interacting degrees of freedom evolving across multiple spatial and temporal scales. Efficient surrogate models therefore require computational architectures capable of preserving local interaction structures while maintaining favorable scaling behavior with increasing system complexity. Finally, uncertainty propagation remains insufficiently understood within current learning-based simulators. Existing architectures typically estimate prediction accuracy indirectly through empirical validation rather than explicitly monitoring the evolution of uncertainty throughout recursive inference. Consequently, prediction confidence often deteriorates long before substantial trajectory divergence becomes observable. Taken together, these challenges indicate that future scientific learning architectures must move beyond purely statistical optimization and instead incorporate explicit computational mechanisms capable of preserving physical structure throughout the entire prediction process. This observation constitutes the primary motivation behind the PHYSGEN framework proposed in this work. \] 
